% Generated draft: requires mathematical and source review.
\documentclass[11pt]{article}
\usepackage[margin=25mm]{geometry}
\usepackage{fontspec}
\setmainfont{Latin Modern Roman}
\usepackage{amsmath,amssymb,mathtools,unicode-math}
\setmathfont{Latin Modern Math}
\usepackage{enumitem,needspace,xurl,hyperref}
\setlength{\parindent}{0pt}
\setlength{\parskip}{5pt}
\setlength{\emergencystretch}{3em}
\newcommand{\N}{\mathbb{N}}
\newcommand{\Z}{\mathbb{Z}}
\newcommand{\Q}{\mathbb{Q}}
\newcommand{\R}{\mathbb{R}}
\newcommand{\C}{\mathbb{C}}
\newcommand{\F}{\mathbb{F}}
\newcommand{\A}{\mathbb{A}}
\newcommand{\PP}{\mathbb{P}}
\newcommand{\E}{\mathbb{E}}
\newcommand{\eps}{\varepsilon}
\newcommand{\dd}{\,\mathrm{d}}
\begin{document}
\section*{UnsolvedMath ID: 20002718}
\textbf{Source problem code:} catalogue entry
\subsection{Definitions and mathematical statement}
D. Armstrong's Problem 4.4 asks for the structure of the assignment from conjugacy classes of a group to Ehrenborg functions of Cayley intervals, and whether it preserves a Hopf-algebra structure. Let \(G\) be generated by a finite conjugation-invariant set \(T\). To fix the undirected Cayley-graph convention, use word length in \(T\cup T^{-1}\), equivalently assume \(T=T^{-1}\) after symmetrization. Write \(\ell\) for this length and set
\[
 x\le_Tg\iff\ell(x)+\ell(x^{-1}g)=\ell(g),\qquad
 P_g=[1,g].
\]
These are finite graded intervals, with rank \(\ell(g)\). Conjugation preserves length and intervals. Translation identifies \([x,g]\) with \(P_{x^{-1}g}\), and the complement map \(u\mapsto u^{-1}g\) gives the relevant order reversal. Thus all subintervals are self-dual. The symmetrization convention and the exact sense of order are not specified in the extracted question.

The printed Ehrenborg sum has missing endpoints and unbounded zero-rank repetitions. Use the following explicit reconstruction. For rank \(d>0\), let \(f_\alpha(P)\) count strict chains from minimum to maximum with successive rank jumps the composition \(\alpha=(\alpha_1,\ldots,\alpha_k)\) of \(d\). Define
\[
 M_\alpha=\sum_{1\le i_1<\cdots<i_k}
       x_{i_1}^{\alpha_1}\cdots x_{i_k}^{\alpha_k},\qquad
 F(P)=\sum_{\alpha\vDash d}f_\alpha(P)M_\alpha,
\]
and set \(F(P)=1\) for the singleton interval. This fixes the usual finite-rank quasisymmetric-function convention without divergent zero-step multiplicities.

The source says that self-duality makes \(F(P)\) symmetric, citing [49]. Global self-duality alone is insufficient: it controls reversal of flag compositions, not every permutation of their parts. The aid reports a rank-four counterexample to that broad assertion and uses the stronger local self-duality of all intervals to obtain \(F(P_g)\in\mathrm{Sym}\). Thus the intended conjugacy-invariant map is
\[
 [g]\longmapsto F(P_g).
\]
The original question asks what algebraic structure underlies this map; it does not specify a product on conjugacy classes.

\textbf{Coalgebra construction reported in the aid.} Over \(\Q\), let \(V_G\) have basis \(e_{[g]}\) indexed by conjugacy classes, graded by \(\ell(g)\). Define
\[
 \delta(e_{[g]})=\sum_{x\le_Tg}
              e_{[x]}\otimes e_{[x^{-1}g]},\qquad
 \varepsilon(e_{[g]})=\begin{cases}1,&g=1,\\0,&g\ne1.\end{cases}
\]
The sum is over actual elements \(x\), so repeated pairs of conjugacy classes retain their multiplicities. The aid reports that it is independent of the representative \(g\), graded and coassociative, with the counit shown, because the two iterated coproducts enumerate reduced three-factorizations. In the usual symmetric-function coproduct,
\[
 \Delta F(P_g)=\sum_{x\le_Tg}F(P_x)\otimes F(P_{x^{-1}g}).
\]
Hence \(e_{[g]}\mapsto F(P_g)\) is a graded coalgebra morphism to \(\mathrm{Sym}\). This alone supplies no canonical multiplication on \(V_G\); group multiplication is not a well-defined graded multiplication of conjugacy classes.

The aid then closes the Cayley intervals under taking intervals and external Cartesian products, forming the associated graded connected incidence Hopf algebra of isomorphism classes. Its product is external poset product, and coproduct is the sum over cuts of an interval. In that completion,
\[
 F(P\times Q)=F(P)F(Q)
\]
and the interval coproduct identity makes \(F\) a Hopf morphism to \(\mathrm{Sym}\). This completion is distinct from asserting a canonical Hopf product on the original fixed-group conjugacy-class space. The source does not identify a unique desired universal completion or further representation-theoretic characterization.

The appended section ``Free Probability'' is background: it credits Voiculescu with the subject, mentions operator-algebra, large-matrix and symmetric-group-representation applications, and credits Speicher with type A noncrossing-partition combinatorics and M\"obius inversion for transforms, referring to survey [47]. It contains no additional question within this record.
\subsection{Short English statement}
The problem asks whether conjugacy classes mapped to Cayley-interval Ehrenborg functions carry compatible algebraic structure. The aid repairs the chain definition and gives a coalgebra morphism with multiplicities. External poset products provide an incidence-Hopf completion, while a canonical product on the fixed group's conjugacy classes remains unspecified. Local self-duality, rather than global self-duality alone, is the symmetry hypothesis used.
\subsection{Sources}
\begin{itemize}
\item Local UnsolvedMath registry, permanent ID 20002718, catalogue entry; canonical \texttt{aim-probability-notes.json}, notes[159], reconstructed-unverified statement status.
\item AIM \emph{Braid groups, clusters and free probability}, Problem 4.4, D. Armstrong, and the appended section 5 background; \url{https://aimath.org/WWN/braidgroups/braidgroups.pdf}. The supplied verification says the repairs are reconstructions and the exact source record remains unchanged.
\item Source reference [49] for the Ehrenborg/symmetry discussion and survey [47] for free probability; full bibliographic entries are absent. Background names Voiculescu and Roland Speicher are retained.
\item Machine-generated research attempt 1, partial result, attempt status unknown, reported rigor proved and low novelty confidence; registry status partially solved. Its coalgebra construction, Hopf completion and self-duality counterexample are not peer reviewed and require independent verification.
\end{itemize}
\end{document}
